第一章已经把群、Lie 群、Lie 代数、表示与不可约表示作为一般数学对象建立起来。把这些对象用于时空对称时，需要始终区分两个物理层次：一方面，Lorentz/Poincar\'e 群本身规定“允许做哪些时空变换以及这些变换怎样复合”；另一方面，一个具体的表示规定同一时空变换怎样作用在坐标四矢量、局域场分量或量子态上。后面出现的矩阵维数不同，并不意味着存在不同的 Lorentz 变换，而只是同一抽象对称作用在不同对象上的实现不同。

对粒子态还需要再强调一次“不可约”的物理意义。若 Hilbert 空间中某个子空间在全部 Poincar\'e 变换下都保持封闭，那么对称性不会把它与其它子空间混合；继续分解直到不能再分，就得到不可约表示。质量、自旋或螺旋度之所以能够成为单粒子种类的标签，正因为它们在这样的最小不变子空间内取固定值。Poincar\'e 代数的 Casimir 不变量在不可约表示上取定值，从而直接给出质量以及自旋或螺旋度标签。


连续群的关键不只是“可以取无穷小参数”，还在于两个无穷小变换的复合必须仍落在同一个群里。取 Hilbert 空间中的厄米生成元 $T_a,T_b$，分别定义
\begin{equation*}
 U_a(\epsilon)=\ee^{-\ii\epsilon T_a/\hbar},
 \qquad
 U_b(\delta)=\ee^{-\ii\delta T_b/\hbar}.
\end{equation*}
若依次沿 $a$ 方向前进、沿 $b$ 方向前进，再沿两条方向逆向返回，得到群交换子
\begin{equation}
\boxed{
 U_a(\epsilon)U_b(\delta)U_a(-\epsilon)U_b(-\delta)
 =\id-\frac{\epsilon\delta}{\hbar^2}[T_a,T_b]
 +\order(\epsilon^2\delta,\epsilon\delta^2).}
\label{eq:group-commutator-bch-dp8}
\end{equation}
这个乘积在 $\epsilon,\delta\to0$ 时仍趋近单位元，因此它的最低阶偏离必须仍属于单位元处的切空间，也就是生成元所张成的 Lie 代数。于是存在与参数无关的结构常数 $f_{ab}{}^c$，使
\begin{equation}
\boxed{
 [T_a,T_b]=\ii\hbar f_{ab}{}^cT_c.}
\label{eq:lie-algebra-closure-dp68}
\end{equation}
这条闭合关系把有限群乘法在恒等元附近的局部结构压缩成生成元的对易代数；Lorentz、Poincar\'e、$SU(2)$ 与 $SU(3)$ 的生成元关系都属于这一共同结构。BCH 展开会把这种闭合性逐阶落实到群乘法系数。

\begin{note}[两种表示中的“酉”]
物理量子态的内积必须在时空对称下保持，因此 Poincar\'e 群在物理 Hilbert 空间上由酉算符表示。局域场分量承担的任务不同：它们必须使场方程在 Lorentz 变换后保持同样的张量形式，所以只要求分量按有限维 Lorentz 表示线性混合。Lorentz 群是非紧群，其非平凡有限维表示一般不能同时保持普通欧几里得分量内积；因此有限维多分量局域场的 Lorentz 变换矩阵通常不是普通意义下的酉矩阵；二分量自旋 $1/2$ 表示中的场分量称为旋量。这与物理态空间上的时间演化和 Poincar\'e 对称必须酉实现并不矛盾。
\end{note}
